\section{Discussion}
The current evidence suggests that the most immediately useful aspect of the system is not automatic route optimization but unified visibility. Even without turning on enforcement modes, a platform team can already observe who spent money, which application violated a declared zone policy, which candidates failed a quality gate, and whether an application bypassed the gateway entirely. That emphasis on observability and governed change aligns more closely with FinOps and AI risk-management practice than with pure benchmark chasing \cite{finopsIntroduction,finopsFramework,finopsGovernance,nistAIRMF,iso42001}.

The completed E2 campaign sharpens that point. Once stale backends and stale API-key secrets were repaired, the limiting factor was no longer the controller path but upstream provider throttling. For this paper, that observation is best read as a protocol warning about live managed-endpoint evaluation rather than as a comparative routing result.

The completed E8 campaign reveals a different weakness. On the current cluster shape, the experiment became scheduling-capacity bound before there was evidence of controller saturation. The central lesson is therefore not that the controller added a large deterministic overhead, but that cluster shape and pod-density ceilings can dominate the measured envelope long before the reconciliation path itself becomes the bottleneck.

A design goal is that governance actuations should be coupled to adequate quality evidence. For the residency-enforcement path this coupling is now enforced: when evidence checking is enabled, a reroute is actuated only on a fresh candidate-safe verdict from a matching quality gate, and otherwise it is withheld and escalated to a human-approved change request (Section~\ref{sec:results}). We demonstrated the escalation branch live and cover the actuation branch with a unit test. The coupling is not yet applied on every actuation path --- budget-driven fallback, in particular, still acts on spend phase alone --- so the invariant holds for residency enforcement but is not yet a global property of the control plane; extending it uniformly is future work.

The shadow-AI follow-up also exposed a practical systems lesson. The eBPF signal itself was present for both OpenAI and Anthropic rogue traffic, but the evidentiary bridge into the operator depended on the \texttt{shadow-egress} handoff remaining populated from current Tetragon exports. In other words, the hard problem was not endpoint recognition but operationalizing the bridge as a continuously refreshed telemetry path.

The manuscript also has to be explicit about control-plane security limits. The latest implementation no longer leaves the approval and evidence paths entirely unprotected: when trusted keys are configured, \texttt{AIChangeRequest} decisions and opt-in evidence ConfigMaps are checked by validating webhooks, and the AKS results show that unsigned or tampered writes are denied. That is a real strengthening of the control plane and gives the paper a clearer identity than a pure future-work note. At the same time, the safe interpretation remains bounded. The measured property is admission-time integrity for two governance-sensitive object classes, not universal tamper evidence under privilege compromise, because a sufficiently privileged cluster administrator could still alter webhook configuration, trusted keys, or the underlying nodes. The paper should therefore claim measured signed-governance admission, not complete control-plane immutability.
